\subsubsection{Memorijski mapovan uređaj i optimizacije}\label{gdb-mmap-volatile}

Neke greške se vide tek kada program komunicira sa memorijom čiju vrednost može da promeni neko van običnog toka C programa: hardver, kernel, drugi proces ili signal handler. Primer ispod simulira uređaj kroz datoteku \texttt{dev.bin} mapiranu u memoriju.

\codefile{main.c}{c}

Program čeka da uređaj više nije zauzet, zatim upisuje podatak i komandu. Ako kompajler zaključi da se polje \texttt{isBusy} ne menja unutar petlje, optimizovana verzija može da ukloni ponovno čitanje tog polja. Popravljena verzija označava pokazivač kao \texttt{volatile}:

\codefile{main-fixed.c}{c}

Prevedimo obe verzije:

\begin{minted}{bash}
$ make
\end{minted}

Ako je datoteka uređaja u stanju ``zauzet'', pokretanje loše verzije može izgledati kao običan beskonačan rad:

\begin{minted}{bash}
$ ./make_device_ready
$ ./sample
\end{minted}

U drugom terminalu prikačimo \texttt{gdb} na proces ili ga pokrenimo direktno:

\begin{minted}{text}
(gdb) break SendCommand
(gdb) run
(gdb) disassemble /m SendCommand
(gdb) info registers
\end{minted}

Komanda \texttt{disassemble /m} spaja izvorni kod i instrukcije. Kod optimizovanog programa to je često važnije od linijskog koraka kroz C kod: linije mogu biti preuređene, spojene ili uklonjene.

Zatim uporedimo popravljenu verziju:

\begin{minted}{text}
(gdb) file ./sample-fixed
(gdb) break SendCommand
(gdb) run
(gdb) disassemble /m SendCommand
\end{minted}

U popravljanoj verziji očekujemo da se polje \texttt{isBusy} ponovo čita iz memorije tokom petlje. To ne znači da \texttt{volatile} rešava sinhronizaciju između niti. Ovde se koristi za komunikaciju sa memorijom čija promena nije vidljiva kroz obična pravila optimizacije C programa.

Realna upotreba ovog obrasca je debagovanje drajvera u korisničkom prostoru, rada sa \texttt{mmap} interfejsom, deljene memorije i grešaka koje se pojavljuju samo u optimizovanim buildovima. U takvim slučajevima treba gledati i C promenljive i generisane instrukcije.
